% Page 069 — page imprimée 65
% La famille Roussel du Villaret au Château d'Ayres

\begin{center}
{\Large\bfseries La famille Roussel}

\vspace{1em}

{\Large\bfseries du Villaret au Château d’Ayres}
\end{center}

\medskip

\textbf{Il était une fois à Ayres}, hameau dominant la vallée de la Jonte, non loin de Meyrueis, un
château. Bien moins imposant mais plus ensoleillé et accueillant que celui de Roquedols, il
avait été construit au 7° siècle. Après bien des vicissitudes, une communauté monastique s’y
installe en 1025 puis l’abandonne au début du 16° siècle. Le château est acheté alors par une
famille protestante de Montpellier, les « de Galtier de Montauran », qui le fortifient.
Christophe de Galtier, chef de guerre énergique et violent, dirigera au nom du roi de Navarre,
la prise et le sac de La Canourgue en 1591 au cours duquel l’église fut incendiée.

\medskip

Lors de la guerre des camisards, en 1702, la troupe de Castanet suivi de la Rose, son fidèle
compagnon, vint camper au Château d’Ayres. \textit{« On fit une assemblée, relate Philippe
Chambon dans « Itinéraires Protestants en Languedoc », présidée par Castanet où se
rendirent les habitants de la cité, mais aussi, des soldats de la garnison, des officiers et même
des prêtres. A l’issue de ce culte une vingtaine de protestants de Meyrueis furent arrêtés.
Réaction immédiate des Camisards qui capturent une quinzaine de soldats et des officiers du
régiment de Cordes. Après négociation chacun relâcha ses otages et l’affaire se conclut par un souper, au château
d’Ayres, réunissant les chefs révoltés et le major du régiment qui assura n’avoir procédé aux
arrestations qu’à l’instigation des prêtres, qui prétendaient que la trêve signée par Cavaler et
Villars ne concernait pas les autres camisards. »}

\medskip

Plus tard le sénéchal d’Anduze, sur les ordres de Richelieu, démantela les deux tours ainsi que
les remparts du château et finalement après la Révocation de l’Edit de Nantes en 1685, Galtier
de Montauran fit amende honorable, abandonna le protestantisme pour sauver ses biens : le
château et son moulin, le moulin d’Ayres, ainsi que de nombreuses terres.
Mal lui en prit car les camisards l’incendièrent. Ses descendants « dilapidèrent l’héritage ».

\medskip

Déjà en 1736 Jean de Galtier de Montauran avait cédé le moulin du Château (Le Moulin
d’Ayres) à son meunier, \textit{« baillé à titre de loyer perpétuel et emphytéose, à Jean Rozier
originaire du Moulin d’Aumières ».}

\medskip

Mais c’est surtout le dernier Charles de Galtier de Montauran qui, joueur invétéré, fut dans
l’obligation de vendre, par morceaux successifs, une bonne partie des terres du château afin
de payer ses dettes de jeux.

\medskip

C’est ainsi qu’en janvier 1839, il doit régler 3000F au sieur Auradou, boulanger, 3.000F aussi
au sieur Agulhon et enfin 900F à un certain Jean Ribes. C’étaient de grosses sommes. Enfin le
château avec le bosquet est vendu au 19° siècle à une famille Nogaret qui voulait le
réhabiliter mais se décide assez rapidement à le revendre. Deux de leurs enfants reposent au
cimetière de Meyrueis.
